
%(BEGIN_QUESTION)
% Copyright 2011, Tony R. Kuphaldt, released under the Creative Commons Attribution License (v 1.0)
% This means you may do almost anything with this work of mine, so long as you give me proper credit

Les og lag en disposisjon av delkapittelet ``RF Link Budget'' i delen ``Radio Systems'' i kapittelet ``Wireless Instrumentation'' i læreboken {\it Lessons In Industrial Instrumentation}. Noter sidetall for viktige illustrasjoner, fotografier, ligninger, tabeller og andre relevante detaljer. Forbered deg på å diskutere begreper og eksempler fra lesingen med lærer og medstudenter.

\underbar{file i00547}
%(END_QUESTION)





%(BEGIN_ANSWER)


%(END_ANSWER)





%(BEGIN_NOTES)

RF-effekten som mottas av en radiomottaker må være høy nok til at signalet ikke drukner i omgivelsesstøy. Vi må ``budsjettere'' alle gevinster og tap i et radiosystem (fra sender til mottaker, inkludert alle former for effekttap) for å sikre god dataintegritet fra ende til ende.

$$P_{rx} = P_{tx} + G_{total} + L_{total}$$

$$P_{tx} = N_f + N_{rx} + S - (G_{total} + L_{total})$$

\noindent
Der,

$P_{rx}$ = Signaleffekt tilgjengelig ved mottakerinngang, i dBm

$P_{tx}$ = Senders utgangssignaleffekt, i dBm

$G_{total}$ = Summen av alle gevinster (forsterkere, antenneretningsvirkning, osv.), positiv dB-verdi

$L_{total}$ = Summen av alle tap (kabler, filtre, path loss, fading, osv.), negativ dB-verdi

$N_f$ = Støygulv, i dBm

$N_{rx}$ = Støytall for mottaker, i dBm

$S$ = Ønsket signal/støy-margin, i dB

\vskip 10pt

Radioprodusenter kapsler ofte inn støygulv, støytall og signal/støy-forhold i én størrelse kalt {\it receiver sensitivity} (minste signalnivå ved mottaker for god dataintegritet).

$$P_{tx} = P_{rx} - (G_{total} + L_{total})$$

\noindent
Der,

$P_{tx}$ = Minste nødvendige utgangssignaleffekt fra sender, i dBm

$P_{rx}$ = Mottakerfølsomhet (minste mottatte signalnivå), i dBm

$G_{total}$ = Summen av alle gevinster (forsterkere, antenneretningsvirkning, osv.), positiv dB-verdi

$L_{total}$ = Summen av alle tap (kabler, filtre, path loss, fading, osv.), negativ dB-verdi

\vskip 10pt

Mottatt signaleffekt korrelerer direkte med datafeil (Bit Error Rate, BER). Jo høyere mottatt effekt, desto sjeldnere korrupte bit.

\vskip 10pt

Path loss er det uunngåelige tapet i signaleffekt som følger av invers-kvadrat-loven når RF-energi sprer seg i rommet etter å ha forlatt sendeantennen. Beregning av path loss i fri sikt:

$$L_p = -20 \log \left(4 \pi D \over \lambda\right)$$

\noindent
Der,

$L_p$ = Path loss, negativ dB-verdi

$D$ = Avstand mellom sende- og mottakerantenne

$\lambda$ = Bølgelengde for RF-signalet, i samme enhet som $D$

\vskip 10pt

Fading oppstår når RF-bølger interfererer destruktivt og svekker mottatt effekt. RF-energi reflektert fra ledende objekter og som ankommer mottaker i motfase, vil svekke signalet. Fading er vanskelig å forutsi kvantitativt, og derfor legger man ofte inn en ``fade margin'' på 20 dB til 30 dB i linkbudsjettet.

\vskip 10pt

Linkbudsjettberegninger er kun estimater. En stedstest anbefales vanligvis for å bestemme nøyaktig sendeeffektbehov.








\vskip 10pt

\filbreak

\vskip 20pt \vbox{\hrule \hbox{\strut \vrule{} {\bf Forslag til sokratisk diskusjon} \vrule} \hrule}

\begin{itemize}
\item{} Forklar hva {\it receiver sensitivity} er, og hvordan dette brukes i linkbudsjettberegninger.
\item{} Forklar mekanismen bak {\it path loss} for en elektromagnetisk bølge.
\item{} Forklar mekanismen bak {\it fade loss} for en elektromagnetisk bølge.
\item{} Øker, minker eller forblir path loss den samme når avstanden mellom antennene øker?
\item{} Øker, minker eller forblir path loss den samme når avstanden mellom antennene minker?
\item{} Øker, minker eller forblir path loss den samme når signalfrekvensen øker?
\item{} Øker, minker eller forblir path loss den samme når signalfrekvensen minker?
\end{itemize}


%INDEX% Leseoppgave: Lessons In Industrial Instrumentation, trådløs instrumentering (RF link budget)

%(END_NOTES)

